% filepath: /home/kluo/Documents/repos/kluo/mathphy/homework/2025/hw1_v2.tex
\documentclass[11pt]{article}

\usepackage[margin=1in]{geometry} 
\usepackage{amsmath,amsthm,amssymb,mathtools}
\usepackage{graphicx, multicol, array}
\usepackage{ctex} 
\usepackage{fancyhdr}
\usepackage{fontspec}
\usepackage{fourier-orns}
\usepackage{xcolor}
\usepackage{enumitem}
\usepackage[most]{tcolorbox}

% Fix headheight warning
\setlength{\headheight}{15.12pt}
\addtolength{\topmargin}{-3.12pt}

% Chinese punctuation handling
\let\oldjuhao=。
\catcode`\。=\active
\newcommand{。}{\ifmmode\text{\oldjuhao}\else\oldjuhao\fi}

\let\olddouhao=，
\catcode`\，=\active
\newcommand{，}{\ifmmode\text{\olddouhao}\else\olddouhao\fi}

% Mathematical symbols
\newcommand{\N}{\mathbb{N}}
\newcommand{\Z}{\mathbb{Z}}
\newcommand{\R}{\mathbb{R}}
\newcommand{\C}{\mathbb{C}}
\newcommand{\ii}{\mathrm{i}}
% Remove conflicting Re and Im definitions since they're already defined
% \DeclareMathOperator{\Re}{Re}
% \DeclareMathOperator{\Im}{Im}

% Problem environment with better styling
\theoremstyle{definition}
\newtheorem{problem}{题目}[section]

% Colors for better visual hierarchy
\definecolor{headercolor}{RGB}{25,25,112}
\definecolor{problemcolor}{RGB}{70,130,180}

% Enhanced header rule
\renewcommand{\headrule}{%
\vspace{-8pt}\hrulefill
\raisebox{-2.0pt}{\quad\decofourleft\decotwo\decofourright\quad}\hrulefill}

% Date configuration
\newcommand{\duedate}{\zhdate{2024/10/08}}
\pagestyle{fancy}

% Header and footer setup
\fancyhf{} % clear all fields
\fancyhead[L]{\color{headercolor}\textbf{数学物理方法}}
\zhnumsetup{style={Traditional,Financial}}
\fancyhead[R]{\color{headercolor}\textbf{课后习题\zhnumber{1}}}

\fancyfoot[L]{\small 截止日期: \duedate}
\fancyfoot[C]{\thepage}
\fancyfoot[R]{\small 任课教师: 罗凯}

% Title page styling
\title{
    \vspace{-1cm}
    {\Huge\color{headercolor}\textbf{数学物理方法}}\\
    \vspace{0.5cm}
    {\Large 课后习题一}
}
\author{
    \textbf{任课教师:} 罗凯\\
    \textbf{截止日期:} \duedate
}
\date{}

\begin{document}

% Custom enumerate styles
\setlist[enumerate,1]{label=(\arabic*), leftmargin=*, itemsep=0.5em}
\setlist[enumerate,2]{label=(\arabic{enumi}.\arabic*), leftmargin=*, itemsep=0.3em}
\setlist[itemize]{leftmargin=*, itemsep=0.3em}

\maketitle
\thispagestyle{fancy}

% Simplified problem box styling (compatible with basic tcolorbox)
\tcbset{
    problembox/.style={
        colback=blue!5!white,
        colframe=problemcolor,
        fonttitle=\bfseries,
        coltitle=white,
        title style={colback=problemcolor}
    }
}

% \section{复数基础}

\begin{tcolorbox}[problembox, title=题目 1]
写出下列复数的实部、虚部、模和辐角。
\begin{enumerate}
    \item $1 + \ii \sqrt{3}$
    \item $\displaystyle\frac{1-\ii}{1+\ii}$
    \item $e^{z}$（其中 $z = x + \ii y$）
    \item $e^{1+\ii}$
    \item $e^{\ii \varphi(x)}$，其中 $\varphi(x)$ 是实数 $x$ 的实函数
\end{enumerate}
\end{tcolorbox}

\begin{tcolorbox}[problembox, title=题目 2]
证明以下规律：
\begin{enumerate}
    \item 复数的加法和乘法满足结合律，即
    \begin{itemize}
        \item $(z_1 + z_2 ) + z_3 = z_1 +( z_2 + z_3)$
        \item $(z_1 \cdot z_2) \cdot z_3 = z_1 \cdot (z_2 \cdot z_3)$
    \end{itemize}
    \item 复数乘积的共轭等于复数共轭的乘积，即
    \[(z_1 \cdot z_2)^* = z_1^* z_2^*\]
\end{enumerate}
\end{tcolorbox}

\section{复平面几何}

\begin{tcolorbox}[problembox, title=题目 3]
找出复平面上满足方程
\[|z - \ii a| = \lambda |z + \ii a|, \quad \lambda > 0\]
的所有点 $z = x + \ii y$。请分 $\lambda$ 的三种情况分别绘制图形：
\begin{enumerate}
    \item $\lambda < 1$
    \item $\lambda > 1$ 
    \item $\lambda = 1$
\end{enumerate}
\end{tcolorbox}

\begin{tcolorbox}[problembox, title=题目 4]
若 $|z| = 1$，$a, b$ 为任意复数，试证明
\[\left| \frac{az + b}{\bar{b}z + \bar{a}} \right| = 1\]
\end{tcolorbox}

% \section{复数的幂与根}

\begin{tcolorbox}[problembox, title=题目 5]
试用 $\cos\varphi$、$\sin\varphi$ 表示 $\cos 4\varphi$。
\end{tcolorbox}

\begin{tcolorbox}[problembox, title=题目 6]
求下列方程的根，并在复平面上画出它们的位置。
\begin{enumerate}
    \item $z^4 + 1 = 0$
    \item $z^2 + 2z\cos\lambda + 1 = 0$，其中 $0 < \lambda < \pi$
\end{enumerate}
\end{tcolorbox}

% \section{复函数}

\begin{tcolorbox}[problembox, title=题目 7]
验证
\[\tan^{-1}(z) = \frac{\ii}{2}[\ln(1-\ii z) - \ln(1+\ii z)]\]
\end{tcolorbox}

\begin{tcolorbox}[problembox, title=题目 8]
试证明极坐标下的柯西-黎曼方程：
\[\left\{\begin{aligned}
\frac{\partial u}{\partial \rho} &= \frac{1}{\rho} \frac{\partial v}{\partial \varphi}\\
\frac{1}{\rho} \frac{\partial u}{\partial \varphi} &= -\frac{\partial v}{\partial \rho}
\end{aligned}\right.\]
\end{tcolorbox}

\begin{tcolorbox}[problembox, title=题目 9]
已知解析函数 $f(z)$ 的实部 $u(x, y)$ 或虚部 $v(x, y)$，求该解析函数。
\begin{enumerate}
    \item $u = e^x \sin y$
    \item $u = x^2 - y^2 + xy$，且 $f(0) = 0$
\end{enumerate}
\end{tcolorbox}

\vfill
\begin{center}
\textcolor{gray}{\hrulefill \quad 作业要求 \quad \hrulefill}\\
\vspace{0.3cm}
\small
请工整书写，步骤完整，逻辑清晰。\\
如有疑问，请及时联系任课教师。
\end{center}

\end{document}